\section{Conclusiones}
A partir de la implementación física del mecanismo y del sistema de adquisición
de datos, se identificaron diversas limitaciones de carácter mecánico y
electrónico que comprometieron la precisión de las mediciones y la
confiabilidad del sistema en general. A continuación se presentan las
principales observaciones y las mejoras propuestas para una versión futura del
prototipo.

\subsection*{Juego mecánico en las articulaciones}

Se evidenció una cantidad considerable de juego mecánico en dos puntos
críticos del mecanismo: la unión entre la tuerca y la barra roscada, y la
articulación $C$. Este juego introdujo movimientos no deseados durante la
operación, lo cual afectó directamente la calidad de las mediciones al generar
desplazamientos no contemplados en el modelo cinemático. Como solución, la
siguiente versión del prototipo incorporará una barra adicional que garantice
el alineamiento de las tuercas y que, al crear un segundo punto de apoyo,
elimine efectivamente el juego presente en la configuración actual.

\subsection*{Gestión del cableado entre sensores y microcontrolador}

Las conexiones eléctricas de largo recorrido entre los sensores y el
microcontrolador fueron realizadas mediante cables UTP, lo cual generó
múltiples dificultades durante el montaje y redujo la confiabilidad del
sistema. La longitud y flexibilidad de estos cables propiciaron desconexiones
intermitentes y complicaron el proceso de ensamblaje. En la siguiente versión
se reemplazarán por cables de cobre de sección adecuada, dirigidos por debajo
del montaje de forma ordenada, con el fin de mejorar tanto la robustez de las
conexiones como la mantenibilidad del sistema.

\subsection*{Tolerancias en piezas fabricadas por manufactura aditiva}

Las piezas producidas mediante impresión 3D presentaron reiterados problemas
de tolerancia dimensional, los cuales requirieron intervención manual mediante
lijado, taladrado y otros procesos de acabado para lograr el ensamblaje
correcto. Esto incrementó el tiempo de montaje y comprometió la repetibilidad
del proceso de fabricación. La siguiente versión adoptará tolerancias más
holgadas en el diseño o, en su defecto, revisará el paradigma de diseño para
evitar ajustes de alta precisión que sean difíciles de alcanzar con los
parámetros de impresión disponibles.

\subsection*{Confiabilidad del sistema de conexiones electrónicas}

Las protoboards empleadas para el prototipado electrónico comenzaron a
presentar fallas de contacto intermitentes debido a su naturaleza desmontable,
lo que introdujo fuentes de error difíciles de diagnosticar durante la
operación. Dado que esta característica es inherente a las protoboards, la
siguiente versión del sistema migrará hacia una PCB dedicada, lo cual
eliminará este tipo de fallas, reducirá el ruido eléctrico y proporcionará una
plataforma electrónica más robusta y reproducible.
